\chapter{Random Integers}

\begin{center}\itshape
One sees, from this Essay, that the theory of probabilities is basically just common sense reduced to calculus.\\
\upshape --- Pierre-Simon Laplace, \emph{Essai philosophique sur les probabilités} (1814)
\end{center}

On 30 January 1812 Carl Friedrich Gauss wrote to Pierre-Simon Laplace, the master of probability, with a question he could not
settle. Take a number $x$ between $0$ and $1$ at random and expand it as a continued fraction,
$x=1/(a_1+1/(a_2+1/(a_3+\cdots)))$. Strip off the first $n$ partial quotients and look at what is left, the number
$x_n=1/(a_{n+1}+1/(a_{n+2}+\cdots))$. What is the probability that $x_n<t$? Gauss had found the answer twelve years earlier,
in a diary entry of 25 October 1800: as $n$ grows, the probability tends to $\log(1+t)/\log2$. What he could not do, he told
Laplace, was to estimate how fast. His proof of the limit has never been found.

The question waited more than a century. In 1928 the Russian mathematician Rodion Kuzmin proved Gauss's claim with an error
that decreases like $q^{\sqrt n}$, and a year later Paul Lévy, independently, got the better error $q^n$ with $q\approx0.7$. The
best possible $q$ was found only in 1974 by Eduard Wirsing: $0.3036630\ldots$, a constant with no known closed form. Along
the way the problem had created a new field, the ergodic theory of numbers, in which a number is a point that moves and its
digits are what an observer records.

Gauss's question is a question about random numbers that have arithmetic inside. This chapter is full of them. What is the
chance that two integers have no common factor, that a number has no square factor, that a tree in an orchard hides
another? None of these questions makes sense at first sight, because there is no uniform probability on the integers. They
make perfect sense as limits, and the answers are products over the primes: a random integer behaves as if its residues
modulo different primes were independent coins.

\begin{think}
Choose two large integers at random. Which is more likely: that they have no common factor, or that neither of them is
divisible by a square larger than $1$? Make a guess, then think about one prime at a time.
\end{think}

\section{Densities}

A set $A$ of positive integers has \emph{density} $d(A)$ if $\frac1N\#(A\cap\{1,\dots,N\})\to d(A)$. For pairs or tuples of
integers we use the same definition with $\{1,\dots,N\}^k$. ``A random integer has property $P$ with probability $c$'' means that
the integers with property $P$ have density $c$. The first problems explain why this language is safe, and the next ones show
the main technique: local conditions at finitely many primes are periodic, hence easy, and the remaining primes are controlled
by a tail estimate.

\begin{exercise}[$\star$]\label{pb:nouniform}
\emph{There is no uniform integer.}
\begin{enumerate}[label=(\alph*)]
\item Show that there is no countably additive probability measure on $\N^*$ that gives probability $\frac1m$ to every residue
class modulo every $m$.
\item Show that density is finitely additive on sets that have a density, and translation invariant.
\item Find two sets $A$ and $B$ that both have density $\frac12$ while $A\cap B$ has no density.
\end{enumerate}
\end{exercise}

\begin{exercise}[$\star$]\label{pb:valuation}
\emph{How divisible is a random integer?} Let $p$ be prime and $n$ uniform in $\{1,\dots,N\}$.
\begin{enumerate}[label=(\alph*)]
\item Show that $P(v_p(n)=k)\to(1-\frac1p)p^{-k}$ for every $k\ge0$, and that $\mathbb E[v_p(n)]\to\frac1{p-1}$.
\item Show that for distinct primes $p_1,\dots,p_r$ the valuations $v_{p_1}(n),\dots,v_{p_r}(n)$ are asymptotically independent.
\item Show that $\mathbb E\bigl[2^{v_2(n)}\bigr]=\frac12\log_2N+O(1)$, while $P\bigl(2^{v_2(n)}\le2\bigr)\to\frac34$. Explain the
paradox.
\end{enumerate}
\end{exercise}

\begin{exercise}[$\star\star$]\label{pb:consecsqfree}
\emph{Squarefree neighbours.}
\begin{enumerate}[label=(\alph*)]
\item Show that the density of the integers $n$ such that $n$ and $n+1$ are both squarefree is
$\prod_p\bigl(1-\frac2{p^2}\bigr)=0.3226\ldots$
\item Show that $n,n+1,n+2$ are all squarefree with density $\prod_p\bigl(1-\frac3{p^2}\bigr)=0.1254\ldots$, and that no four
consecutive integers are all squarefree.
\item Show that there are arbitrarily long runs of consecutive integers none of which is squarefree.
\end{enumerate}
\end{exercise}

\begin{exercise}[$\star\star$]\label{pb:pairwise}
\emph{Pairwise coprime.} Let $k\ge2$.
\begin{enumerate}[label=(\alph*)]
\item Show that $k$ random integers are pairwise coprime with probability
\[
\prod_p\Bigl(1-\frac1p\Bigr)^{k-1}\Bigl(1+\frac{k-1}p\Bigr).
\]
\item For $k=3$ this is $0.2867\ldots$, larger than $(6/\pi^2)^3=0.2247\ldots$, the value if the three events
``$\gcd(a,b)=1$'', ``$\gcd(b,c)=1$'', ``$\gcd(a,c)=1$'' were independent. Explain the sign of the difference prime by prime.
\end{enumerate}
\end{exercise}

\begin{exercise}[$\star$]\label{pb:parity}
\emph{A random fraction.} Choose $a,b$ at random and write $\frac ab=\frac{a'}{b'}$ in lowest terms.
\begin{enumerate}[label=(\alph*)]
\item Show that $p\mid b'$ with probability $\frac1{p+1}$.
\item Deduce that the three parity patterns odd/odd, odd/even, even/odd of $\frac{a'}{b'}$ are equally likely.
\item Show that the events $p\mid b'$, for distinct primes $p$, are asymptotically independent.
\end{enumerate}
\end{exercise}

\begin{exercise}[$\star\star$]\label{pb:nsq1sqfree}
\emph{Squarefree values of $n^2+1$.}
\begin{enumerate}[label=(\alph*)]
\item Show that the density of the $n$ such that $n^2+1$ is squarefree is
\[
\prod_{p\equiv1\ (4)}\Bigl(1-\frac2{p^2}\Bigr)=0.8948\ldots
\]
\item Why is this so much larger than the density $\frac6{\pi^2}=0.6079\ldots$ of all squarefree integers?
\end{enumerate}
\end{exercise}

\begin{exercise}[$\star\star\star$]\label{pb:twosqfree}
\emph{Two squarefree summands.} Show that every integer $n\ge2$ is the sum of two squarefree positive integers. Prove it for
$n\ge14000$ by counting, and check the smaller $n$ by computer.
\end{exercise}

\begin{exercise}[$\star\star$]\label{pb:phiovern}
\emph{The law of $\varphi(n)/n$.}
\begin{enumerate}[label=(\alph*)]
\item Show that for every integer $k\ge1$,
\[
\frac1N\sum_{n\le N}\Bigl(\frac{\varphi(n)}n\Bigr)^k\longrightarrow\prod_p\Bigl(1-\frac1p+\frac1p\Bigl(1-\frac1p\Bigr)^k\Bigr).
\]
\item Deduce that $\frac1N\sum_{n\le N}g\bigl(\frac{\varphi(n)}n\bigr)$ has a limit for every continuous $g:[0,1]\to\R$: the values
$\varphi(n)/n$ have a limiting distribution.
\item Show that at least half of all integers, asymptotically, have $\varphi(n)\le\frac n2$, although the mean of
$\varphi(n)/n$ is $\frac6{\pi^2}>\frac12$.
\end{enumerate}
\end{exercise}

\section{Common factors and lattice points}

Two random integers have a gcd, which is $1$ more than half of the time and occasionally large. Seen in the plane, a pair
$(a,b)$ is a lattice point, and $\gcd(a,b)=1$ means that nothing stands between it and the origin.

\begin{exercise}[$\star\star$]\label{pb:gcdmean}
\emph{The law of the gcd.} Let $a,b$ be independent and uniform in $\{1,\dots,N\}$.
\begin{enumerate}[label=(\alph*)]
\item Show that $P(\gcd(a,b)=g)\to\frac6{\pi^2g^2}$ for every $g\ge1$.
\item Show that $\mathbb E[\gcd(a,b)]=\frac6{\pi^2}\log N+O(1)$.
\item Show that $\mathbb E[\log\gcd(a,b)]\to-\frac{\zeta'(2)}{\zeta(2)}=0.5699\ldots$
\item Show that $\mathbb E[\operatorname{lcm}(a,b)]\sim\frac{\zeta(3)}{4\zeta(2)}N^2=0.1827\ldots N^2$.
\end{enumerate}
\end{exercise}

\begin{exercise}[$\star$]\label{pb:fibcoprime}
\emph{Coprime Fibonacci numbers.} Let $a,b$ be uniform in $\{1,\dots,N\}$ and $F_n$ the Fibonacci numbers.
\begin{enumerate}[label=(\alph*)]
\item Show that $P(\gcd(F_a,F_b)=1)\to\frac{15}{2\pi^2}=0.7599\ldots$
\item Show that $\mathbb E[\log\gcd(F_a,F_b)]=\frac{6\log\varphi}{\pi^2}\log N+O(1)$, where $\varphi$ is the golden ratio.
\end{enumerate}
\end{exercise}

\begin{exercise}[$\star\star$]\label{pb:visible}
\emph{Visible points.} A point $(a,b)\in\Z^2\setminus\{0\}$ is \emph{visible} from the origin if no other lattice point lies on the
segment joining them.
\begin{enumerate}[label=(\alph*)]
\item Show that $(a,b)$ is visible if and only if $\gcd(a,b)=1$.
\item Let $V(R)$ be the number of visible points in the disc of radius $R$. Show that $V(R)=\frac6\pi R^2+O(R\log R)$.
\item Show that the visible points are equidistributed in angle: the proportion of them in a sector of angle $\theta$ tends to
$\frac\theta{2\pi}$.
\end{enumerate}
\end{exercise}

\begin{figure}[ht]
\centering
\fig{ch06-visible}
\caption{The lattice points of the disc of radius $11$: visible from the origin (filled) or hidden behind a nearer point (rings).
A proportion close to $\frac6{\pi^2}\approx61\%$ are visible (Problem~\ref{pb:visible}). The red sight line, towards $(11,4)$,
passes between the trees: in Pólya's orchard (Problem~\ref{pb:orchard}) it escapes when the trunks are thin enough.}
\label{fig:visible}
\end{figure}

\begin{exercise}[$\star\star$]\label{pb:orchard}
\emph{Pólya's orchard (1918).} Let $R\ge2$ be an integer. A tree stands at every lattice point of the disc of radius $R$ except
the origin, and each trunk is a disc of radius $r$. You stand at the origin.
\begin{enumerate}[label=(\alph*)]
\item (Minkowski.) Show that a closed rectangle centred at $0$, of area at least $4$, contains a lattice point other than $0$.
\item Show that if $r\ge\frac1R$, every line of sight from the origin meets a trunk.
\item Show that if $r<\frac1{\sqrt{R^2+1}}$, you can see out of the orchard.
\end{enumerate}
\end{exercise}

\begin{exercise}[$\star\star$]\label{pb:farey}
\emph{Farey fractions.} The Farey sequence $\mathcal F_N$ lists the fractions $\frac hk$ in $[0,1]$ in lowest terms with $1\le k\le N$,
in increasing order.
\begin{enumerate}[label=(\alph*)]
\item Show that $|\mathcal F_N|=1+\sum_{k\le N}\varphi(k)=\frac3{\pi^2}N^2+O(N\log N)$.
\item Show that two consecutive fractions $\frac hk<\frac{h'}{k'}$ of $\mathcal F_N$ satisfy $h'k-hk'=1$ and $k+k'>N$.
\item Deduce that $\sum\frac1{kk'}=1$, the sum running over consecutive pairs, and that the gaps of $\mathcal F_N$ all lie between
$\frac1{N(N-1)}$ and $\frac1N$.
\end{enumerate}
\end{exercise}

\section{Euclid's algorithm at random}

Euclid's algorithm on $(a,b)$ computes the continued fraction of $\frac ab$, and its quotients are the partial quotients. Gauss's
question of the opening is about these quotients for a random real number. The section ends with a full proof of his limit law,
with an explicit exponential error, and its consequence for the growth of the denominators.

\begin{exercise}[$\star\star$]\label{pb:euclidgame}
\emph{Euclid's game (Cole and Davie, 1969).} Two players start from a pair $(a,b)$ of positive integers. A move replaces the larger
number by itself minus a positive multiple of the smaller, as long as the result is $\ge0$. The player who produces a $0$ wins.
\begin{enumerate}[label=(\alph*)]
\item Let $\varphi=\frac{1+\sqrt5}2$. Show that, for $a<b$, the player to move wins if and only if $\frac ba>\varphi$.
\item Deduce that, for random $a<b$, the first player wins with probability tending to $\frac1\varphi=0.618\ldots$
\end{enumerate}
\end{exercise}

\begin{exercise}[$\star$]\label{pb:gaussmeasure}
\emph{The Gauss measure.} Let $T(x)=\{1/x\}$ on $]0,1[$ (fractional part), so that $a_{n+1}(x)=\lfloor1/T^nx\rfloor$ is the
$(n+1)$-st partial quotient of $x$, and let $\mu$ be the probability measure with density $\frac1{(1+x)\log2}$ on $[0,1]$.
\begin{enumerate}[label=(\alph*)]
\item Show that $\mu(T^{-1}[0,t[)=\mu([0,t[)$ for every $t\in[0,1]$: $\mu$ is invariant under $T$.
\item Deduce that, if $x$ has law $\mu$, then $P(a_n=k)=\log_2\frac{(k+1)^2}{k(k+2)}$ for every $n\ge1$ and $k\ge1$.
\item If $x$ is uniform instead, compute $P(a_1=k)$, and compare with (b) for $k=1$.
\end{enumerate}
\end{exercise}

\begin{exercise}[Journey]\label{pb:gausskuzmin}
\emph{The Gauss--Kuzmin theorem.} Let $x$ be uniform in $]0,1[$, $m_n(t)=P(T^nx<t)$ and $f_n=m_n'$, so $f_0=1$.
\begin{enumerate}[label=(\alph*)]
\item Show that $m_{n+1}(t)=\sum_{k\ge1}\bigl(m_n(\frac1k)-m_n(\frac1{k+t})\bigr)$, and deduce
$f_{n+1}(t)=\sum_{k\ge1}\frac1{(k+t)^2}f_n\bigl(\frac1{k+t}\bigr)$, with every $f_n$ continuously differentiable.
\item Let $\phi_n(t)=(1+t)f_n(t)$. Show that $\phi_{n+1}(t)=\sum_kP_k(t)\phi_n(u_k)$ with $u_k=\frac1{k+t}$ and
$P_k(t)=\frac{1+t}{(k+t)(k+t+1)}$, and that $\sum_kP_k(t)=1$.
\item Show that $\sup|\phi_{n+1}'|\le c\,\sup|\phi_n'|$ with
\[
c=\frac12+2\sum_{k\ge2}\frac1{k^3(k+1)}+\sum_{k\ge2}\frac{k-1}{k(k+1)^2}=2\zeta(3)-\frac32=0.9041\ldots
\]
\item Deduce that $\Bigl|f_n(t)-\frac1{(1+t)\log2}\Bigr|\le\frac{c^n}{1+t}$, and conclude: $|m_n(t)-\log_2(1+t)|\le c^n\log(1+t)$.
\item Deduce that $P(a_n=k)\to\log_2\frac{(k+1)^2}{k(k+2)}$ for $x$ uniform.
\end{enumerate}
\end{exercise}

\begin{figure}[ht]
\centering
\fig{ch06-gausskuzmin}
\caption{The density $f_n$ of $T^nx$ for $x$ uniform in $]0,1[$: $f_0=1$, $f_1(t)=\sum_k(k+t)^{-2}$, and $f_2$, computed numerically,
already within $0.06$ of Gauss's density $\frac1{(1+t)\log2}$ (dashed). Problem~\ref{pb:gausskuzmin} proves that the error shrinks
at least like $0.905^n$; Wirsing showed the true rate is $0.3037^n$.}
\label{fig:gausskuzmin}
\end{figure}

\begin{exercise}[$\star\star$]\label{pb:levy}
\emph{How fast the denominators grow.} Let $x$ be uniform in $]0,1[$ and $\frac{p_n}{q_n}$ its convergents.
\begin{enumerate}[label=(\alph*)]
\item Show that $|q_nx-p_n|=x\,Tx\cdots T^nx$, and that this number lies between $\frac1{2q_{n+1}}$ and $\frac1{q_{n+1}}$.
\item Show that $\int_0^1\frac{-\log t}{1+t}\,dt=\frac{\pi^2}{12}$.
\item Using Problem~\ref{pb:gausskuzmin}, prove that
\[
\mathbb E[\log q_n]=\frac{\pi^2}{12\log2}\,n+O(1),\qquad\frac{\pi^2}{12\log2}=1.18656\ldots
\]
\end{enumerate}
\end{exercise}

\section{Digits, residues and walks}

Integers are also strings of digits and collections of residues. The last problems look at a random integer through these
two windows, and end with a random walk that mixes in a number of steps only polylogarithmic in the size of its space.

\begin{exercise}[$\star\star$]\label{pb:benford}
\emph{Leading digits.} Let $A_d$ be the set of positive integers whose decimal expansion starts with the digit $d$.
\begin{enumerate}[label=(\alph*)]
\item Show that $A_1$ has no density: the proportion of $n\le N$ in $A_1$ oscillates between $\frac19$ and $\frac59$.
\item Show that $A_d$ has \emph{logarithmic density} $\log_{10}(1+\frac1d)$:
\[
\frac1{\log N}\sum_{n\le N,\ n\in A_d}\frac1n\longrightarrow\log_{10}\Bigl(1+\frac1d\Bigr).
\]
\end{enumerate}
\end{exercise}

\begin{figure}[ht]
\centering
\fig{ch06-benford}
\caption{The proportion of $n\le N$ whose first digit is $1$, for $N$ up to $10^5$ on a logarithmic scale. It climbs to almost
$\frac59$ at $N=2\cdot10^k-1$ and falls back to $\frac19$ at $N=10^k-1$, forever (Problem~\ref{pb:benford}). Averaged with weights
$\frac1n$ it settles at Benford's value $\log_{10}2=0.301$.}
\label{fig:benford}
\end{figure}

\begin{exercise}[$\star\star$]\label{pb:digitsclt}
\emph{Binary digits.} Let $s(n)$ be the number of $1$'s in the binary expansion of $n$, and $n$ uniform in $\{0,\dots,N-1\}$.
\begin{enumerate}[label=(\alph*)]
\item If $N=2^k$, show that $s(n)$ is binomial with parameters $k$ and $\frac12$.
\item For general $N$, show that $\bigl(s(n)-\frac12\log_2N\bigr)/\sqrt{\frac14\log_2N}$ converges in law to a standard Gaussian.
\end{enumerate}
\end{exercise}

\begin{exercise}[$\star\star$]\label{pb:squaresmodpk}
\emph{Squares modulo $p^k$.} Let $S(m)$ be the number of squares in $\Z/m\Z$ ($0$ included).
\begin{enumerate}[label=(\alph*)]
\item For an odd prime $p$, show that $S(p^k)=1+\sum_{0\le j<k/2}\frac{p-1}2p^{k-2j-1}$, and that $S(p^k)/p^k\to\frac p{2(p+1)}$.
\item Show that $S(2^k)/2^k\to\frac16$.
\item Show that $S$ is multiplicative, and deduce that the proportion of squares modulo $m$ can be made as small as we like.
What does this say about an integer that is a square modulo every $m$?
\end{enumerate}
\end{exercise}

\begin{exercise}[$\star$]\label{pb:hashing}
\emph{Universal hashing (Carter and Wegman, 1979).} Let $p$ be prime, $a,b$ independent and uniform in $\Z/p\Z$, and
$h(x)=\bigl((ax+b)\bmod p\bigr)\bmod m$ for keys $x\in\{0,\dots,p-1\}$, where $2\le m\le p$.
\begin{enumerate}[label=(\alph*)]
\item Show that for $x\ne y$ the pair $(ax+b,ay+b)$ is uniform on $(\Z/p\Z)^2$.
\item Deduce that $P(h(x)=h(y))\le\frac1m+\frac1p$ for $x\ne y$.
\item Store $n$ keys in $m=n^2$ buckets with such an $h$. Show that with probability at least $\frac12-\frac{n^2}{2p}$ no two keys
collide.
\end{enumerate}
\end{exercise}

\begin{exercise}[$\star\star\star$]\label{pb:cdg}
\emph{A fast random walk modulo $p$ (Chung, Diaconis and Graham, 1987).} Let $p$ be an odd prime, $X_0=0$ and
$X_{n+1}=2X_n+\varepsilon_n\bmod p$, where the $\varepsilon_n$ are independent and uniform in $\{-1,0,1\}$. Let $d_n$ be the total
variation distance between the law of $X_n$ and the uniform law.
\begin{enumerate}[label=(\alph*)]
\item Show that $d_n\ge1-3^n/p$, so that $\log_3p$ steps are necessary.
\item Show that $\mathbb E\bigl[e^{2i\pi tX_n/p}\bigr]=\prod_{j<n}\frac{1+2\cos(2\pi2^jt/p)}3$.
\item Let $L=\lceil\log_2p\rceil+1$. Show that among any $L$ consecutive $j$, at least one has $\cos(2\pi2^jt/p)\le0$ (for
$t\not\equiv0$).
\item Deduce that $d_n\le\frac12\sqrt p\,3^{-\lfloor n/L\rfloor}$, so that $O(\log^2p)$ steps suffice.
\end{enumerate}
\end{exercise}

\section{Hints}

\begin{hint}{pb:nouniform}
(a) What would the probability of a single integer be? (c) Let $B$ agree with the even numbers on some long blocks and with the
odd numbers on others.
\end{hint}

\begin{hint}{pb:valuation}
(a) $\#\{n\le N:v_p(n)\ge k\}=\lfloor N/p^k\rfloor$. (c) The mean is dominated by the rare $n$ divisible by a large power of $2$.
\end{hint}

\begin{hint}{pb:consecsqfree}
(a) For the primes $p\le y$ the condition is periodic modulo $\prod_{p\le y}p^2$; the primes $p>y$ exclude at most $2N/y+O(\sqrt N)$
integers. (c) Chinese remainders.
\end{hint}

\begin{hint}{pb:pairwise}
(a) Locally at $p$: at most one of the $k$ integers is divisible by $p$. (b) Compare the two local factors.
\end{hint}

\begin{hint}{pb:parity}
(a) $p\mid b'$ exactly when $v_p(b)>v_p(a)$.
\end{hint}

\begin{hint}{pb:nsq1sqfree}
Count the roots of $x^2\equiv-1$ modulo $p^2$. For the tail, $p^2\mid n^2+1$ with $n\le N$ forces $p\le N$.
\end{hint}

\begin{hint}{pb:twosqfree}
Inclusion--exclusion over the primes up to $13$, which leaves an error of at most $4^6$ terms, and a tail estimate for the
primes from $17$ on.
\end{hint}

\begin{hint}{pb:phiovern}
(a) $(\varphi(n)/n)^k$ is multiplicative and close to $1$ at large primes: use Wintner's theorem (Tome~5, Theorem~2.2).
(b) Weierstrass.
\end{hint}

\begin{hint}{pb:gcdmean}
(b) $\gcd(a,b)=\sum_{d\mid a,\ d\mid b}\varphi(d)$. (d) Pull out the gcd and use Möbius for the coprime part.
\end{hint}

\begin{hint}{pb:fibcoprime}
$\gcd(F_a,F_b)=F_{\gcd(a,b)}$ (Tome~0, Exercise~1.7).
\end{hint}

\begin{hint}{pb:visible}
(b) Every non-zero lattice point is a multiple $d\cdot v$ of a unique visible point $v$; invert with Möbius.
\end{hint}

\begin{hint}{pb:orchard}
(a) Halve the rectangle and look at its translates by $\Z^2$. (b) A rectangle of length $2R$ and width $2r$ around the line
of sight; a lattice point $z$ with $|z|^2<R^2+1$ has $|z|\le R$. (c) Look towards $(R,1)$.
\end{hint}

\begin{hint}{pb:farey}
(b) Solve $kx-hy=1$ with $N-k<y\le N$, and show that $\frac xy$ is the successor of $\frac hk$.
\end{hint}

\begin{hint}{pb:euclidgame}
(a) Induction on $a+b$. When $b>2a$ the mover can choose between two positions, and exactly one of them has ratio below
$\varphi$.
\end{hint}

\begin{hint}{pb:gaussmeasure}
(a) $T^{-1}[0,t[=\bigcup_k\,]\frac1{k+t},\frac1k]$, and the series telescopes.
\end{hint}

\begin{hint}{pb:gausskuzmin}
(c) Differentiate the recursion of (b). Since $\sum P_k'=0$, the first sum is $\sum P_k'(t)\bigl(\phi_n(u_k)-\phi_n(u_1)\bigr)$.
(d) $\frac1{\log2}$ is a weighted average of $\phi_n$.
\end{hint}

\begin{hint}{pb:levy}
(a) $x=\frac{p_n+p_{n-1}T^nx}{q_n+q_{n-1}T^nx}$ (Tome~1, Proposition~8.1). (c) Take expectations in $\log q_{n+1}=-\sum_{j\le n}\log T^jx+O(1)$.
\end{hint}

\begin{hint}{pb:benford}
(b) The block $[d\cdot10^j,(d+1)10^j[$ contributes about $\log(1+\frac1d)$ to the sum.
\end{hint}

\begin{hint}{pb:digitsclt}
(b) Split $\{0,\dots,N-1\}$ according to the binary expansion of $N$ into dyadic blocks, on each of which $s(n)$ is a shifted
binomial.
\end{hint}

\begin{hint}{pb:squaresmodpk}
Write a non-zero square as $p^{2j}u$ with $u$ a unit; the units that are squares form a subgroup of index $2$ ($4$ for $p=2$).
\end{hint}

\begin{hint}{pb:hashing}
(c) Expected number of colliding pairs, and Markov's inequality.
\end{hint}

\begin{hint}{pb:cdg}
(c) If $|\theta|<\frac14$ then doubling $\theta$ modulo $1$ is just doubling. (d) Plancherel on $\Z/p\Z$.
\end{hint}

\section{Solutions}

\begin{solution}{pb:nouniform}
\idea A uniform probability would give every single integer the same mass, which must be $0$, and countable additivity then
kills everything. Densities avoid the contradiction by giving up countable additivity.

\writeup (a) Suppose $P$ is such a measure. A singleton $\{n\}$ lies in the class of $n$ modulo $m$, so $P(\{n\})\le\frac1m$ for
every $m$, hence $P(\{n\})=0$, and $P(\N^*)=\sum_nP(\{n\})=0$, a contradiction.

(b) If $A,B$ are disjoint with densities, $\#((A\cup B)\cap[1,N])=\#(A\cap[1,N])+\#(B\cap[1,N])$, so $A\cup B$ has density
$d(A)+d(B)$. If $A$ has a density, $\#((A+h)\cap[1,N])$ differs from $\#(A\cap[1,N])$ by at most $|h|$, which is $o(N)$.

(c) Let $E$ be the even numbers and $I_j=[4^j,2\cdot4^j[$. Let $B$ be the set of even numbers in $\bigcup_jI_j$ together with the
odd numbers outside it. In any interval, $B$ contains one of every two consecutive integers, so $B$ has density $\frac12$, and
so does $E$. But $E\cap B$ is the set of even numbers of $\bigcup I_j$. The blocks $I_j$ make up a proportion of $[1,N]$ that is
about $\frac23$ at $N=2\cdot4^j$ and about $\frac13$ at $N=4^{j+1}$, so the proportion of $E\cap B$ oscillates between about
$\frac13$ and $\frac16$, and has no limit.

\tools Tome~3, Chapter~1 (additivity of probabilities).
\leads Densities are finitely additive ``probabilities'' on the sets that have one. By the Hahn--Banach theorem (Part~II) they
extend to all subsets as invariant means, but no such extension can be written down explicitly.
\end{solution}

\begin{solution}{pb:valuation}
\idea Divisibility by $p^k$ happens exactly once every $p^k$ integers, so the valuation is geometric in the limit. Distinct
primes do not interact, by the Chinese remainder theorem.

\writeup (a) The number of $n\le N$ with $p^k\mid n$ is $\lfloor N/p^k\rfloor=N/p^k+O(1)$, so
$P(v_p(n)=k)=\frac1N(\lfloor N/p^k\rfloor-\lfloor N/p^{k+1}\rfloor)\to p^{-k}-p^{-k-1}$. For the mean,
$\mathbb E[v_p(n)]=\sum_{k\ge1}P(v_p(n)\ge k)=\frac1N\sum_{k\ge1}\lfloor N/p^k\rfloor$, which lies between
$\sum_{k\le\log_pN}(p^{-k}-\frac1N)$ and $\sum_kp^{-k}=\frac1{p-1}$, and tends to $\frac1{p-1}$.

(b) For integers $k_i\ge0$, the condition $p_i^{k_i}\mid n$ for all $i$ is $m\mid n$ with $m=\prod p_i^{k_i}$, which has
probability $\frac1N\lfloor N/m\rfloor\to\frac1m=\prod_ip_i^{-k_i}$. So the joint tail probabilities $P(v_{p_i}(n)\ge k_i\ \forall i)$
converge to products, and by inclusion--exclusion so do the probabilities $P(v_{p_i}(n)=k_i\ \forall i)$.

(c) $2^{v_2(n)}\le2$ means $4\nmid n$, of probability $\to\frac34$. For the mean, the number of $n\le N$ with $v_2(n)=k$ is
$\frac N{2^{k+1}}+\theta_k$ with $|\theta_k|\le1$, and it is $0$ for $2^k>N$. So
\[
\sum_{n\le N}2^{v_2(n)}=\sum_{k\le\log_2N}2^k\Bigl(\frac N{2^{k+1}}+\theta_k\Bigr)=\frac N2\bigl(\lfloor\log_2N\rfloor+1\bigr)+O(N),
\]
since $\sum_{k\le\log_2N}2^k\le2N$. Divide by $N$. There is no paradox: the mean is carried by rare values. The $n$ with
$v_2(n)=k$ have probability about $2^{-k-1}$ and contribute $2^k\cdot2^{-k-1}=\frac12$ each to the mean, and there are
$\log_2N$ such levels.

\tools Tome~1, Chapter~1 (valuations); Tome~2, Chapter~2 (Chinese remainders).
\leads In the language of Tome~1, Chapter~12, a random integer is a sample of Haar measure on $\Z_p$, simultaneously for all $p$:
that is the meaning of (b). The quantity $2^{v_2(n)}$ is the ``largest power of $2$ dividing $n$'', whose average governs the
cost of the binary gcd algorithm.
\end{solution}

\begin{solution}{pb:consecsqfree}
\idea Squarefreeness is a condition at every prime. For finitely many primes it is periodic and its density is a finite
product; for the large primes, being divisible by $p^2$ is so rare that the total effect vanishes.

\writeup (a) Fix $y\ge2$, let $Q=\prod_{p\le y}p^2$, and let $S_y$ be the set of $n$ such that neither $n$ nor $n+1$ is divisible by
$p^2$ for any $p\le y$. Membership in $S_y$ depends only on $n\bmod Q$, and by the Chinese remainder theorem the number of good
residues modulo $Q$ is $\prod_{p\le y}(p^2-2)$: modulo $p^2$ we must avoid $n\equiv0$ and $n\equiv-1$. So
$\#(S_y\cap[1,N])=N\prod_{p\le y}(1-\frac2{p^2})+O(Q)$.

Let $G$ be the set of $n$ with $n$ and $n+1$ squarefree. Then $G\subset S_y$, and an $n\le N$ in $S_y\setminus G$ has $p^2\mid n$ or
$p^2\mid n+1$ for some prime $p>y$ with $p^2\le N+1$. Their number is at most
\[
\sum_{y<p\le\sqrt{N+1}}2\Bigl(\frac N{p^2}+1\Bigr)\le\frac{2N}y+2\sqrt{N+1},
\]
using $\sum_{p>y}\frac1{p^2}\le\sum_{m>y}\frac1{m(m-1)}\le\frac1y$. Dividing by $N$ and letting $N\to\infty$,
\[
\prod_{p\le y}\Bigl(1-\frac2{p^2}\Bigr)-\frac2y\le\liminf\frac{\#(G\cap[1,N])}N\le\limsup\frac{\#(G\cap[1,N])}N\le\prod_{p\le y}\Bigl(1-\frac2{p^2}\Bigr).
\]
Now let $y\to\infty$; the product converges since $\sum\frac2{p^2}<\infty$.

(b) The same argument, with three forbidden residues $0,-1,-2$ modulo each $p^2$, gives $\prod_p(1-\frac3{p^2})$; for $p=2$ the
three residues $0,3,2$ modulo $4$ are distinct, so the factor is $\frac14$, and the product is not $0$. Among four consecutive
integers one is divisible by $4$.

(c) Given $L$, choose distinct primes $p_1,\dots,p_L$ and, by the Chinese remainder theorem, $n$ with $n\equiv-j\pmod{p_j^2}$ for
$j=1,\dots,L$. Then $p_j^2\mid n+j$, and none of $n+1,\dots,n+L$ is squarefree.

\tools Tome~2, Chapter~2 (Chinese remainders); Tome~4, Exercise~10.2 (the density $\frac6{\pi^2}$ of squarefree numbers).
\leads By (c) the first run of $L$ non-squarefree integers starts below $\prod_{j\le L}p_j^2$, and in practice much
earlier: the notebook finds that a run of $4$ first starts at $242$ ($242=2\cdot11^2$, $243=3^5$, $244=4\cdot61$, $245=5\cdot7^2$).
How long the longest run below $x$ is remains open; a heuristic count (each of the $L$ integers needs its own prime square,
of size about $L$) suggests about $\frac{\log x}{2\log\log x}$.
\end{solution}

\begin{solution}{pb:pairwise}
\idea At each prime the $k$ integers behave like independent coins of bias $\frac1p$, and ``pairwise coprime at $p$'' means at
most one coin shows heads. The probability is the product of these local probabilities, by the same truncation as in
Problem~\ref{pb:consecsqfree}.

\writeup (a) For a prime $p$, the probability that at most one of $k$ independent residues modulo $p$ is $0$ is
$(1-\frac1p)^k+k\frac1p(1-\frac1p)^{k-1}=(1-\frac1p)^{k-1}(1+\frac{k-1}p)$. Fix $y$ and let $Q=\prod_{p\le y}p$. The condition ``no
prime $p\le y$ divides two of the $a_i$'' depends only on the $a_i$ modulo $Q$, and by Chinese remainders its density in
$\{1,\dots,N\}^k$ is $c_y=\prod_{p\le y}(1-\frac1p)^{k-1}(1+\frac{k-1}p)$, up to $O(Q/N)$. The $k$-tuples satisfying it but not
pairwise coprime have two entries divisible by a common prime $p>y$; there are at most
$\binom k2\sum_{p>y}N^{k-2}\lfloor N/p\rfloor^2\le\binom k2N^k/y$ of them. Let $N\to\infty$, then $y\to\infty$, as before.

(b) The local factor in (a) for $k=3$ is $(1-\frac1p)^2(1+\frac2p)=1-\frac3{p^2}+\frac2{p^3}$, whereas $(1-\frac1{p^2})^3=1-\frac3{p^2}+\frac3{p^4}-\frac1{p^6}$.
The difference is $\frac2{p^3}-\frac3{p^4}+\frac1{p^6}>0$. The three events are positively correlated at every prime: if $p$
divides $a$ and $b$, the events $p\mid\gcd(b,c)$ and $p\mid\gcd(a,c)$ happen together, so the bad events overlap and the good
event is more likely than independence predicts.

\tools Tome~2, Chapter~2; Tome~3, Chapter~1 (inclusion--exclusion, independence); Tome~5, Exercise~2.3 (the gcd of $k$ integers).
\leads For $k=4$ the probability is $0.1149\ldots$ As $k\to\infty$ it tends to $0$ very fast, because each small prime
$p$ is then very likely to divide at least two of the $k$ integers. The same local reasoning gives the probability that
the $\binom k2$ differences of $k$ random integers are all coprime to $m$.
\end{solution}

\begin{solution}{pb:parity}
\idea Reducing $\frac ab$ removes the common power of $p$, so $p$ stays in the denominator exactly when $b$ has more factors $p$
than $a$. The valuations of $a$ and $b$ are independent geometric variables.

\writeup (a) $p\mid b'$ if and only if $v_p(b)>v_p(a)$. By Problem~\ref{pb:valuation}(a) applied to $a$ and $b$ independently, for
each $J$ the probability that $v_p(b)>v_p(a)$ and $v_p(a)<J$ tends to
\[
\sum_{j<J}\Bigl(1-\frac1p\Bigr)p^{-j}\cdot p^{-j-1},
\]
and the event $v_p(a)\ge J$ has probability at most $p^{-J}$. Letting $J\to\infty$, the probability tends to
$\frac{p-1}{p^2}\cdot\frac1{1-p^{-2}}=\frac1{p+1}$.

(b) By (a) and symmetry, $2\mid b'$ and $2\mid a'$ each have probability $\frac13$, and they cannot both happen. So each of the
three patterns odd/odd, odd/even, even/odd has probability $\frac13$.

(c) For distinct primes $p_1,\dots,p_r$, the event ``$v_{p_i}(b)>v_{p_i}(a)$ and $v_{p_i}(a)<J$ for all $i$'' is a union of
conditions on $(a,b)$ modulo $\prod p_i^{J+1}$, and by Problem~\ref{pb:valuation}(b) its probability tends to the product of the
individual ones; letting $J\to\infty$ as in (a) gives $\prod_i\frac1{p_i+1}$.

\tools Problem~\ref{pb:valuation}; Tome~1, Chapter~1.
\leads Among random fractions in lowest terms the parity classes are equidistributed, a fact that reappears in the action of the
modular group $SL_2(\Z)$ on $\mathbb P^1(\mathbb F_2)$, a set of $3$ points permuted transitively. By (c), the reduced denominator
of a random fraction is divisible by $m$ with probability $\prod_{p\mid m}\frac1{p+1}$ for squarefree $m$.
\end{solution}

\begin{solution}{pb:nsq1sqfree}
\idea The local factor at $p$ is $1-\rho(p^2)/p^2$, where $\rho(p^2)$ counts the solutions of $n^2\equiv-1\pmod{p^2}$. The novelty
is the tail: the primes up to $N$ are now all relevant, but a polynomial of degree $2$ has few roots, and the number of primes
is $o(N)$.

\writeup Let $\rho(m)$ be the number of $x\bmod m$ with $x^2\equiv-1$. For $p=2$, $x^2+1\equiv1,2\pmod4$, so $\rho(4)=0$. For $p\equiv3\pmod4$,
$-1$ is not a square modulo $p$, so $\rho(p^2)=0$. For $p\equiv1\pmod4$ there are two roots modulo $p$, and each lifts uniquely to
a root modulo $p^2$: if $x_0^2\equiv-1\pmod p$, then $(x_0+tp)^2+1\equiv x_0^2+1+2x_0tp\pmod{p^2}$, and since $2x_0\not\equiv0$ there is
exactly one $t$ modulo $p$. So $\rho(p^2)=2$. In all cases $\rho(p^2)\le2$.

As in Problem~\ref{pb:consecsqfree}, the $n$ with $p^2\nmid n^2+1$ for all $p\le y$ have density $\prod_{p\le y}(1-\rho(p^2)/p^2)$, which is
$\prod_{p\le y,\,p\equiv1(4)}(1-\frac2{p^2})$. For the tail: if $p^2\mid n^2+1$ with $n\le N$, then $p^2\le N^2+1<(N+1)^2$, so $p\le N$. The
number of $n\le N$ with $p^2\mid n^2+1$ is at most $\rho(p^2)(\frac N{p^2}+1)\le\frac{2N}{p^2}+2$. Summing over $y<p\le N$ gives at most
$\frac{2N}y+2\pi(N)$, and $\pi(N)=o(N)$ (Tome~5, Chapter~5, or already Chebyshev's bound). The rest is as before.

(b) The density of squarefree numbers is $\prod_p(1-\frac1{p^2})$: every prime removes a proportion $\frac1{p^2}$, and $p=2$ and $p=3$
remove the most, $\frac14$ and $\frac19$. For $n^2+1$, the primes $2$ and $3$ (and all $p\equiv3\pmod4$) remove nothing: $n^2+1$ is never
divisible by $4$ or by $9$. The first prime that matters is $5$, which removes $\frac2{25}$.

\tools Tome~1, Chapter~6 ($-1$ is a square modulo $p$ exactly when $p\equiv1\pmod4$); Tome~1, Chapter~12 (Hensel's lifting);
Problem~\ref{pb:consecsqfree}.
\leads Estermann (1931) proved (a) with an error term. For a polynomial of degree $3$, squarefree values are known to have the
expected density (Hooley, Erdős for $k$-free values), but for degree $\ge4$ the question is open in general: the tail estimate
fails, because the primes up to $N^{d/2}$ matter and there are too many of them. Granville showed in 1998 that the $abc$
conjecture would settle all degrees.
\end{solution}

\begin{solution}{pb:twosqfree}
\idea The expected number of representations $n=a+(n-a)$ with both parts squarefree is about $0.32\,n$. To turn this into a
proof one needs an \emph{explicit} error, so the sieve uses only the primes up to $13$, where inclusion--exclusion has few terms,
and a crude tail bound for the others.

\writeup Let $A(n)$ be the number of $a\in\{1,\dots,n-1\}$ with $a$ and $n-a$ squarefree, $P=2\cdot3\cdot5\cdot7\cdot11\cdot13$, and
$B(n)$ the number of $a$ in the same range such that $p^2\nmid a$ and $p^2\nmid n-a$ for every $p\mid P$. By inclusion--exclusion,
\[
B(n)=\sum_{d\mid P}\sum_{e\mid P}\mu(d)\mu(e)\,\#\{1\le a\le n-1:\ d^2\mid a,\ e^2\mid n-a\}.
\]
The two congruences $a\equiv0\pmod{d^2}$, $a\equiv n\pmod{e^2}$ are compatible if and only if $g^2\mid n$, where $g=\gcd(d,e)$, and
then they define one class modulo $L=\operatorname{lcm}(d^2,e^2)$, which contains $\frac{n-1}L+\theta$ elements of $[1,n-1]$, with
$|\theta|<1$. So $B(n)=(n-1)\Pi(n)+E$ with $|E|\le4^6=4096$ (the number of pairs $(d,e)$), and $\Pi(n)$ the multiplicative sum
\[
\Pi(n)=\sum_{d,e\mid P}\frac{\mu(d)\mu(e)[g^2\mid n]}L=\prod_{p\mid P}\Bigl(1-\frac2{p^2}+\frac{[p^2\mid n]}{p^2}\Bigr)\ge\prod_{p\le13}\Bigl(1-\frac2{p^2}\Bigr)=0.33350\ldots
\]
(the sum factors over primes because the summand is a product of local factors, each coming from $d,e\in\{1,p\}$).

An $a$ counted by $B(n)$ but not by $A(n)$ has $p^2\mid a$ or $p^2\mid n-a$ for some prime $17\le p\le\sqrt n$; there are at most
$2\sum_{p\ge17}\frac n{p^2}<0.0331\,n$ such $a$, since $\sum_{p\ge17}p^{-2}=0.016546\ldots$ (this is the prime zeta value
$\sum_pp^{-2}=0.452247\ldots$ minus its first six terms). Hence
\[
A(n)\ge0.3335(n-1)-4096-0.0331\,n>0.30\,n-4097,
\]
which is positive for $n\ge13657$. For $2\le n<14000$ the notebook finds a representation directly (for instance $2=1+1$, $3=1+2$),
and in fact the least $a$ that works is always very small.

\tools Tome~1, Chapter~4 (Möbius function, multiplicativity); Tome~2, Chapter~2 (Chinese remainders); Problem~\ref{pb:consecsqfree}.
\leads Estermann (1931) showed that the number of representations is $\sim c(n)\,n$ with an explicit singular series
$c(n)=\prod_p(1-\frac2{p^2}+\frac{[p^2\mid n]}{p^2})$, the exact analogue of the Hardy--Littlewood formula for Goldbach's problem,
with squarefree numbers in place of primes. For primes the same conjectural formula is far beyond reach, because primes are
not defined by local conditions alone.
\end{solution}

\begin{solution}{pb:phiovern}
\idea $(\varphi(n)/n)^k$ is a multiplicative function equal to $\sum_{d\mid n}h(d)$ for a function $h$ that vanishes off squarefree
numbers and is $O(k/p)$ at primes. Wintner's theorem then gives the mean, and polynomials approximate every continuous function.

\writeup (a) Let $f_k(n)=(\varphi(n)/n)^k=\prod_{p\mid n}(1-\frac1p)^k$, multiplicative. Define $h$ multiplicative by
$h(p)=(1-\frac1p)^k-1$ and $h(p^j)=0$ for $j\ge2$. Then $\sum_{d\mid n}h(d)=\prod_{p\mid n}(1+h(p))=f_k(n)$. Since $|h(p)|\le\frac kp$,
\[
\sum_d\frac{|h(d)|}d\le\prod_p\Bigl(1+\frac k{p^2}\Bigr)<\infty.
\]
By Wintner's theorem (Tome~5, Theorem~2.2), $f_k$ has mean value $\sum_d\frac{h(d)}d=\prod_p\bigl(1+\frac{h(p)}p\bigr)$, which is the
stated product.

(b) By (a) and linearity, $\frac1N\sum_{n\le N}g(\varphi(n)/n)$ converges for every polynomial $g$. For $g$ continuous and
$\varepsilon>0$, Weierstrass's theorem (Tome~3, Theorem~4.5, via Bernstein polynomials) gives a polynomial $q$ with
$|g-q|\le\varepsilon$ on $[0,1]$. The averages of $g$ and of $q$ differ by at most $\varepsilon$, and those of $q$ form a Cauchy
sequence; so those of $g$ do too, and converge.

(c) Every even $n$ has $\varphi(n)/n\le\frac12$, and the even numbers have density $\frac12$. For $k=1$ the product in (a) is
$\prod_p(1-\frac1{p^2})=\frac6{\pi^2}$. The mean exceeds $\frac12$ because the odd numbers mostly have $\varphi(n)/n$ close to $1$.

\tools Tome~5, Chapter~2 (Wintner's theorem); Tome~3, Chapter~4 (Weierstrass's theorem); Tome~1, Chapter~4.
\leads Schoenberg (1928) proved that the limiting distribution function exists and is continuous, and Erdős (1939) that it is
purely singular: it has no density, although it increases strictly on $[0,1]$. The notebook draws it; the proportion of $n$ with
$\varphi(n)\le n/2$ is about $0.511$: the even numbers, and a few odd numbers with many small prime factors, such as
$105=3\cdot5\cdot7$, for which $\varphi(105)/105=\frac{48}{105}$.
\end{solution}

\begin{solution}{pb:gcdmean}
\idea $\gcd(a,b)=g$ means $g\mid a$, $g\mid b$ and the quotients are coprime: probability $\frac1{g^2}\cdot\frac6{\pi^2}$. The tail
$\sum\frac1{g^2}\cdot g$ diverges logarithmically, which is why the mean grows like $\log N$, while $\log g$ is still integrable.

\writeup (a) $\gcd(a,b)=g$ if and only if $a=ga'$, $b=gb'$ with $a',b'\le N/g$ coprime. By Tome~1, Exercise~4.8, the number of coprime
pairs in $\{1,\dots,M\}^2$ is $\frac6{\pi^2}M^2+O(M\log M)$, so $P(\gcd=g)=\frac1{N^2}\bigl(\frac6{\pi^2}\frac{N^2}{g^2}+O(\frac Ng\log N)\bigr)\to\frac6{\pi^2g^2}$.

(b) By Gauss's identity $\sum_{d\mid m}\varphi(d)=m$ (Tome~1, Proposition~4.5), $\gcd(a,b)=\sum_{d\mid a,\,d\mid b}\varphi(d)$. So
\[
\sum_{a,b\le N}\gcd(a,b)=\sum_{d\le N}\varphi(d)\Bigl\lfloor\frac Nd\Bigr\rfloor^2=N^2\sum_{d\le N}\frac{\varphi(d)}{d^2}+O\Bigl(N\sum_{d\le N}\frac{\varphi(d)}d\Bigr),
\]
because $\lfloor N/d\rfloor^2=N^2/d^2+O(N/d)$. The error is $O(N^2)$. For the main term, $\Phi(x)=\sum_{d\le x}\frac{\varphi(d)}d=\frac6{\pi^2}x+O(\log x)$
(Tome~5, Exercise~2.1), and Abel summation gives
\[
\sum_{d\le N}\frac{\varphi(d)}{d^2}=\frac{\Phi(N)}N+\int_1^N\frac{\Phi(x)}{x^2}\,dx=\frac6{\pi^2}\log N+O(1).
\]
Divide by $N^2$.

(c) $\mathbb E[\log\gcd]=\sum_{g\le N}\log g\,P(\gcd=g)$. Each term tends to $\frac{6\log g}{\pi^2g^2}$, and is bounded by
$\log g\cdot\frac{\lfloor N/g\rfloor^2}{N^2}\le\frac{\log g}{g^2}$, which is summable. By dominated convergence the expectation tends
to $\frac6{\pi^2}\sum_g\frac{\log g}{g^2}=\frac{-\zeta'(2)}{\zeta(2)}$, since $-\zeta'(2)=\sum_g\frac{\log g}{g^2}=0.93754\ldots$

(d) Since $\operatorname{lcm}(a,b)=\frac{ab}{\gcd(a,b)}$, grouping by $g=\gcd(a,b)$,
\[
\sum_{a,b\le N}\operatorname{lcm}(a,b)=\sum_{g\le N}g\,C\Bigl(\frac Ng\Bigr),\qquad C(M)=\sum_{\substack{a',b'\le M\\\gcd(a',b')=1}}a'b'.
\]
By Möbius, $C(M)=\sum_{d\le M}\mu(d)d^2S(M/d)^2$ with $S(x)=\sum_{n\le x}n=\frac{x^2}2+O(x)$; so
$C(M)=\frac{M^4}4\sum_{d\le M}\frac{\mu(d)}{d^2}+O(M^3\log M)=\frac{6}{\pi^2}\frac{M^4}4+O(M^3\log M)$. Then
\[
\sum_{g\le N}g\,C\Bigl(\frac Ng\Bigr)=\frac6{\pi^2}\frac{N^4}4\sum_{g\le N}\frac1{g^3}+O\Bigl(N^3\log N\sum_g\frac1{g^2}\Bigr)=\frac{\zeta(3)}{4\zeta(2)}N^4+O(N^3\log N).
\]
Divide by $N^2$.

\tools Tome~1, Chapter~4 (Exercise~4.8, Proposition~4.5, Möbius inversion); Tome~5, Chapter~2; Tome~3, Chapter~5 (dominated
convergence).
\leads Cesàro studied the mean of the gcd in the 1880s, and the error term in (b) is related to the Dirichlet divisor problem.
The values in (c) and (d) mean that the gcd of two random integers is ``typically'' $e^{0.57}\approx1.77$ but has an infinite
mean in the limit, a heavy tail like the one in Problem~\ref{pb:valuation}(c).
\end{solution}

\begin{solution}{pb:fibcoprime}
\idea Fibonacci numbers inherit their gcds from their indices, and $F_g=1$ only for $g=1,2$.

\writeup (a) By Tome~0, Exercise~1.7, $\gcd(F_a,F_b)=F_{\gcd(a,b)}$. Since $F_1=F_2=1$ and $F_g\ge2$ for $g\ge3$, the gcd is $1$ exactly
when $\gcd(a,b)\in\{1,2\}$, which by Problem~\ref{pb:gcdmean}(a) has probability tending to $\frac6{\pi^2}(1+\frac14)=\frac{15}{2\pi^2}$.

(b) Binet's formula gives $F_g=\frac{\varphi^g-(-1/\varphi)^g}{\sqrt5}$, so $|\log F_g-g\log\varphi+\log\sqrt5|\le C$ for all $g\ge1$,
for some constant $C$. Hence $\mathbb E[\log F_{\gcd(a,b)}]=\log\varphi\cdot\mathbb E[\gcd(a,b)]+O(1)$, and Problem~\ref{pb:gcdmean}(b)
finishes.

\tools Tome~0, Chapter~1 (Fibonacci gcd, Binet's formula); Problem~\ref{pb:gcdmean}.
\leads The same holds for every strong divisibility sequence, for instance $2^n-1$ (Tome~1, Chapter~1), where
$\gcd(2^a-1,2^b-1)=1$ with probability $\frac6{\pi^2}$. Two random Fibonacci numbers are coprime more often than two random integers.
\end{solution}

\begin{solution}{pb:visible}
\idea Each non-zero lattice point is $d$ times a visible point, with $d$ the gcd of its coordinates. So the counting function of all
points is a sum of rescaled counting functions of visible points, and Möbius inversion goes back.

\writeup (a) If $d=\gcd(a,b)>1$, the point $(a/d,b/d)$ lies strictly between $0$ and $(a,b)$. Conversely, a lattice point on the segment
is $t(a,b)$ with $0<t<1$; if $t=\frac uv$ in lowest terms, $v\mid ua$ and $v\mid ub$, so $v\mid\gcd(a,b)$, and $v>1$.

(b) Let $N(R)$ be the number of lattice points in the closed disc of radius $R$. By Tome~5, Exercise~2.4, $|N(R)-\pi R^2|\le C(R+1)$.
Writing each non-zero point as $d\cdot v$ with $v$ visible, $N(R)-1=\sum_{d\ge1}V(R/d)$, and $V(R/d)=0$ for $d>R$. By Möbius inversion
(the sum is finite),
\[
V(R)=\sum_{d\le R}\mu(d)\bigl(N(R/d)-1\bigr)=\pi R^2\sum_{d\le R}\frac{\mu(d)}{d^2}+O\Bigl(\sum_{d\le R}\Bigl(\frac Rd+1\Bigr)\Bigr).
\]
The error is $O(R\log R)$, and $\sum_{d\le R}\frac{\mu(d)}{d^2}=\frac6{\pi^2}+O(\frac1R)$. So $V(R)=\frac6\pi R^2+O(R\log R)$.

(c) Repeat (b) with the sector $\Sigma$ of angle $\theta$ and radius $R$ in place of the disc. The unit squares centred at the lattice
points of $\Sigma$ cover $\Sigma$ minus a band of width $\frac{\sqrt2}2$ along its boundary, and lie in $\Sigma$ plus such a band; the
boundary has length $2R+\theta R$, so the number of lattice points in $\Sigma$ is $\frac\theta2R^2+O(R+1)$. The Möbius inversion is
unchanged (a sector is mapped to the same sector by $v\mapsto dv$), and gives $\frac{6}{\pi^2}\frac\theta2R^2+O(R\log R)$ visible
points. Divide by $V(R)$.

\tools Tome~5, Chapter~2 (Gauss's circle problem); Tome~1, Chapter~4 (Möbius inversion).
\leads The visible points have holes of every size: by Chinese remainders there are $k\times k$ squares of invisible points for every
$k$. In higher dimension the density of visible points is $\frac1{\zeta(n)}$. The gaps between the angles of visible points have a
limiting distribution computed by Boca, Cobeli and Zaharescu (2000), which is not exponential: the angles are equidistributed but
not Poisson.
\end{solution}

\begin{solution}{pb:orchard}
\idea If the trunks are thick, a thin rectangle around any line of sight has area at least $4$ and so contains a tree. If they are
thin, a line of sight with an ``irrational-looking'' slope $\frac1R$ stays far from all lattice points.

\writeup (a) Let $S$ be the rectangle and $K=\frac12S$, of area at least $1$. Suppose first that the area of $K$ is $>1$. The pieces
$K\cap(z+[0,1[^2)$, $z\in\Z^2$, translated back by $-z$ into $[0,1[^2$, have total area $>1$, so two of them overlap: there are
$x\ne y$ in $K$ with $x-y\in\Z^2\setminus\{0\}$. Then $x-y=\frac12(2x)+\frac12(-2y)$ is the midpoint of two points of $S$ ($S$ is
symmetric), so it lies in $S$ (convex). If the area of $K$ is exactly $1$, apply this to $(1+\frac1m)S$ for each $m$: each contains
a non-zero lattice point, and these lie in the bounded set $2S$, so one of them, $z$, lies in $(1+\frac1m)S$ for infinitely many $m$;
since $S$ is closed, $z\in S$.

(b) If $r\ge1$, every direction makes an angle at most $\frac\pi4$ with one of $(\pm1,0),(0,\pm1)$, and the tree there, at distance
at most $\sin\frac\pi4<1$ from the line, blocks it. Let $\frac1R\le r<1$, and let $D$ be a ray from $0$. Take coordinates $(u,v)$ with
$u$ along $D$, and the rectangle $|u|\le R$, $|v|\le r$, of area $4rR\ge4$. By (a) it contains a lattice point $z\ne0$; replacing $z$ by
$-z$ we may take $u\ge0$, and $u>0$ because $|v|\le r<1\le|z|$. Then $|z|^2=u^2+v^2\le R^2+r^2<R^2+1$, and $|z|^2$ is an integer, so
$|z|\le R$: there is a tree at $z$, and its distance $|v|$ to the ray is at most $r$. The trunk meets $D$.

(c) Let $D$ be the ray towards $(R,1)$. The distance from a lattice point $(a,b)$ to the line through $0$ and $(R,1)$ is
$\frac{|a-bR|}{\sqrt{R^2+1}}$. If $a-bR\ne0$ it is at least $\frac1{\sqrt{R^2+1}}>r$, so that trunk misses the line. If $a=bR$, then
$(a,b)=b(R,1)$ with $b\ne0$, and $|(a,b)|\ge\sqrt{R^2+1}>R$: no tree stands there. So no trunk meets $D$.

\tools Tome~0, Chapters~4 and~7 (convexity; determinants as areas); Tome~1, Chapter~1.
\leads This is Problem~239 of Part~VIII in Pólya and Szegő's \emph{Problems and Theorems in Analysis}, after Pólya's 1918 paper.
Part~(a) is the simplest case of Minkowski's theorem, the first theorem of the geometry of numbers, which Tome~12 develops. The
notebook computes the exact critical radius for small $R$: it equals the lower bound $\frac1{\sqrt{R^2+1}}$ every time, so the sight
line towards $(R,1)$ of Figure~\ref{fig:visible} is the best one.
\end{solution}

\begin{solution}{pb:farey}
\idea Neighbours in the Farey sequence span a parallelogram of area $1$ in the lattice. Bézout produces the candidate neighbour,
and a short inequality shows that nothing fits in between.

\writeup (a) The fractions of $\mathcal F_N$ with denominator $k$ are $\frac hk$ with $0\le h\le k$, $\gcd(h,k)=1$: there are $\varphi(k)$ of
them for $k\ge2$ and two ($\frac01,\frac11$) for $k=1$. So $|\mathcal F_N|=1+\sum_{k\le N}\varphi(k)$, and Mertens's estimate
$\sum_{k\le N}\varphi(k)=\frac3{\pi^2}N^2+O(N\log N)$ (Tome~5, Theorem~2.6) concludes.

(b) Let $\frac hk<1$ be in $\mathcal F_N$. Since $\gcd(h,k)=1$, the solutions $(x,y)$ of $kx-hy=1$ form one family $y\equiv y_0\pmod k$,
so there is one with $N-k<y\le N$; then $\gcd(x,y)=1$ and $\frac xy=\frac hk+\frac1{ky}$, and $0<\frac xy\le1$, because $h\le k-1$ gives
$\frac xy\le\frac{k-1}k+\frac1{ky}\le\frac{k-1}k+\frac1k=1$. So $\frac xy\in\mathcal F_N$. Suppose some $\frac{h'}{k'}\in\mathcal F_N$ lies strictly
between. Then $\frac xy-\frac{h'}{k'}\ge\frac1{yk'}$ and $\frac{h'}{k'}-\frac hk\ge\frac1{kk'}$ (positive fractions with these
denominators), so
\[
\frac1{ky}=\frac xy-\frac hk\ge\frac1{yk'}+\frac1{kk'}=\frac{k+y}{kyk'}>\frac N{kyk'}\ge\frac1{ky},
\]
because $k+y>N\ge k'$. This is absurd, so $\frac xy$ is the successor of $\frac hk$: $h'=x$, $k'=y$, $h'k-hk'=1$, and $k+k'>N$.

(c) The gap between consecutive fractions is $\frac{h'}{k'}-\frac hk=\frac1{kk'}$, and the gaps add up to $1-0$. Since $k,k'\le N$, $k\ne k'$
unless $k=k'=1$ (consecutive fractions with the same denominator $k\ge2$ would differ by $\frac1{kk}$, contradicting $h'k-hk'=1$ with
$k\mid1$), so for $N\ge2$ each gap satisfies $\frac1{kk'}\ge\frac1{N(N-1)}$, and $\frac1{kk'}\le\frac1{k\cdot(N+1-k)}\le\frac1N$ by
$k+k'\ge N+1$.

\tools Tome~1, Chapter~1 (Bézout); Tome~5, Chapter~2 (Mertens's estimate for $\sum\varphi$).
\leads Franel (1924) and Landau showed that the Riemann hypothesis is equivalent to the statement that the Farey fractions are
very evenly spread: $\sum_j|f_j-\frac jn|=O(N^{1/2+\varepsilon})$, where $f_1<\dots<f_n$ is $\mathcal F_N$. The mediant
$\frac{h+h'}{k+k'}$ of two neighbours is the first fraction to appear between them, which is how the Stern--Brocot tree of
Chapter~7 is built.
\end{solution}

\begin{solution}{pb:euclidgame}
\idea When $b<2a$ there is only one move, so nothing is decided by the mover. When $b>2a$ the mover can stop at
$(a,a+r)$ or go on to $(r,a)$, and $\varphi$ is exactly the ratio that separates the winning choice from the losing one,
because $\varphi-1=\frac1\varphi$.

\writeup (a) Call a position $(a,b)$ with $a\le b$ \emph{winning} if the player to move can force a win. If $a=b$ the mover wins at once.
We show by induction on $a+b$ that for $a<b$, $(a,b)$ is winning if and only if $b/a>\varphi$ (equality is impossible, $\varphi$ being
irrational). Write $b=qa+r$ with $0\le r<a$, $q\ge1$.

\emph{Case $b/a<\varphi$.} Then $q=1$ (as $\varphi<2$) and $r=b-a>0$, the only move gives $(r,a)$, and $\frac ar>\varphi$ because
$\frac ra=\frac ba-1<\varphi-1=\frac1\varphi$. By induction the opponent, now to move from $(r,a)$, wins.

\emph{Case $b/a>\varphi$.} If $r=0$ the mover wins at once by producing a $0$. Otherwise, if $q=1$, the only move
gives $(r,a)$ with $\frac ar<\varphi$ (since $\frac ra>\frac1\varphi$), a losing position for the opponent. If $q\ge2$, the mover may
leave either $(r,a)$ or $(a,a+r)$. Now $\frac{a+r}a<\varphi$ if and only if $\frac ra<\frac1\varphi$, if and only if $\frac ar>\varphi$; so
exactly one of the two positions has ratio below $\varphi$, and the mover picks it. By induction the opponent loses.

(b) The number of pairs $1\le a<b\le N$ with $b>\varphi a$ is $\sum_{a<N/\varphi}(N-\lfloor\varphi a\rfloor)=\frac{N^2}{2\varphi}+O(N)$,
out of $\binom N2$ pairs; the ratio tends to $\frac1\varphi$.

\tools Tome~1, Chapters~1 and~8 (Euclid's algorithm, continued fractions); Tome~0, Chapter~5 (the golden ratio).
\leads The game follows Euclid's algorithm, and the only real decisions are at partial quotients $\ge2$: the winner is decided by
the first partial quotient of $b/a$ that is at least $2$, by its parity. Cole and Davie's paper appeared in the \emph{Mathematical
Gazette} 53 (1969).
\end{solution}

\begin{solution}{pb:gaussmeasure}
\idea The preimage of $[0,t[$ is a union of intervals, one for each value of the first partial quotient, and their Gauss measures
telescope back to $\log_2(1+t)$.

\writeup (a) $Tx<t$ with $\lfloor1/x\rfloor=k$ means $k\le\frac1x<k+t$, that is $x\in\,]\frac1{k+t},\frac1k]$. So
\[
\mu(T^{-1}[0,t[)=\frac1{\log2}\sum_{k\ge1}\Bigl(\log\Bigl(1+\frac1k\Bigr)-\log\Bigl(1+\frac1{k+t}\Bigr)\Bigr)
=\frac1{\log2}\sum_{k\ge1}\Bigl(\log\frac{k+1}k-\log\frac{k+1+t}{k+t}\Bigr).
\]
The partial sum up to $K$ is $\log(K+1)-\log\frac{K+1+t}{1+t}\to\log(1+t)$, so the value is $\mu([0,t[)$. Two probability measures that
agree on all intervals $[0,t[$ agree on all Borel sets, so $\mu\circ T^{-1}=\mu$.

(b) $a_n(x)=a_1(T^{n-1}x)$, and $T^{n-1}x$ has law $\mu$ by (a) and induction. So
$P(a_n=k)=\mu(]\frac1{k+1},\frac1k])=\log_2\frac{1+1/k}{1+1/(k+1)}=\log_2\frac{(k+1)^2}{k(k+2)}$.

(c) $P(a_1=k)=\frac1k-\frac1{k+1}=\frac1{k(k+1)}$. For $k=1$: $0.5$ for uniform $x$, $\log_2\frac43=0.415$ for the Gauss measure.

\tools Tome~1, Chapter~8; Tome~3, Chapter~5 (laws determined by distribution functions).
\leads $\mathbb E[a_1]=\infty$ under both measures, since the probabilities decay like $\frac1{k^2}$. Khinchin proved that nevertheless
$(a_1\cdots a_n)^{1/n}\to2.6854\ldots$ for almost every $x$, a consequence of the ergodic theorem (Part~II).
\end{solution}

\begin{solution}{pb:gausskuzmin}
\idea The transfer operator replaces $\phi_n$ by an average of its values at the points $u_k$, and these points are squeezed
together: $|u_k'|=(k+t)^{-2}$. Averaging over squeezed points contracts the Lipschitz constant, and a function with tiny
Lipschitz constant is almost constant; the constant is forced by normalisation.

\writeup (a) For irrational $x$, $T^{n+1}x<t$ if and only if $T^nx\in\,]\frac1{k+t},\frac1k]$ for some $k$ (Problem~\ref{pb:gaussmeasure}(a)),
which gives the formula for $m_{n+1}$. Suppose $f_n$ is continuous and bounded by $M_n$ on $[0,1]$. The series
$\sum_k\frac{f_n(1/(k+t))}{(k+t)^2}$ of derivatives of the terms converges uniformly (it is bounded by $M_n\sum k^{-2}$), so $m_{n+1}$ is
differentiable with the stated derivative, continuous and bounded by $M_n\zeta(2)$. The same argument, one derivative higher,
shows by induction that $f_n$ is $C^1$, since $f_0=1$.

(b) $f_n(u_k)=\phi_n(u_k)/(1+u_k)$ and $1+u_k=\frac{k+t+1}{k+t}$, so $\frac{(1+t)f_n(u_k)}{(k+t)^2}=\frac{1+t}{(k+t)(k+t+1)}\phi_n(u_k)=P_k(t)\phi_n(u_k)$.
Also $\sum_kP_k(t)=(1+t)\sum_k\bigl(\frac1{k+t}-\frac1{k+t+1}\bigr)=1$.

(c) Let $\Lambda_n=\sup_{[0,1]}|\phi_n'|$. Differentiating (b) term by term (justified as in (a)),
\[
\phi_{n+1}'(t)=\sum_kP_k'(t)\bigl(\phi_n(u_k)-\phi_n(u_1)\bigr)-\sum_kP_k(t)\frac{\phi_n'(u_k)}{(k+t)^2},
\]
where we subtracted $\phi_n(u_1)\sum_kP_k'=0$. A computation gives
$P_k'(t)=\frac{k^2-k-(1+t)^2}{(k+t)^2(k+t+1)^2}$. For $k=1$ the first term vanishes. For $k\ge2$, $|k^2-k-(1+t)^2|\le k^2$, so
$|P_k'|\le\frac1{(k+1)^2}$, and $|\phi_n(u_k)-\phi_n(u_1)|\le\Lambda_n(u_1-u_k)=\Lambda_n\frac{k-1}{(1+t)(k+t)}\le\Lambda_n\frac{k-1}k$. For the second
sum, $\frac{P_k(t)}{(k+t)^2}=\frac{1+t}{(k+t)^3(k+t+1)}$ is at most $\frac12$ for $k=1$ (it equals $\frac1{(1+t)^2(2+t)}$) and at most
$\frac2{k^3(k+1)}$ for $k\ge2$. Hence $\Lambda_{n+1}\le c\Lambda_n$ with $c$ as stated. With $\frac1{k^3(k+1)}=\frac1{k^3}-\frac1{k^2}+\frac1k-\frac1{k+1}$
and $\frac{k-1}{k(k+1)^2}=\frac1{(k+1)^2}-\frac1{k(k+1)^2}$, $\frac1{k(k+1)^2}=\frac1k-\frac1{k+1}-\frac1{(k+1)^2}$, the sums are
$\sum_{k\ge2}\frac1{k^3(k+1)}=\zeta(3)-\zeta(2)+\frac12$ and $\sum_{k\ge2}\frac{k-1}{k(k+1)^2}=2\zeta(2)-3$, so $c=\frac12+2\zeta(3)-2\zeta(2)+1+2\zeta(2)-3=2\zeta(3)-\frac32=0.9041\ldots$

(d) $\phi_0=1+t$ has $\Lambda_0=1$, so $\Lambda_n\le c^n$ and $\max\phi_n-\min\phi_n\le c^n$. Since $\int_0^1f_n=1$,
$\int_0^1\phi_n(t)\frac{dt}{1+t}=1$: the number $\frac1{\log2}$ is the average of $\phi_n$ for the probability $\frac{dt}{(1+t)\log2}$,
so it lies between $\min\phi_n$ and $\max\phi_n$, and $|\phi_n(t)-\frac1{\log2}|\le c^n$. Divide by $1+t$, and integrate from $0$ to $t$:
$|m_n(t)-\log_2(1+t)|\le c^n\log(1+t)$.

(e) $P(a_{n+1}=k)=P(T^nx\in\,]\frac1{k+1},\frac1k])=m_n(\frac1k)-m_n(\frac1{k+1})$, which by (d) tends to
$\log_2(1+\frac1k)-\log_2(1+\frac1{k+1})$.

\tools Tome~0, Chapter~5 (uniform convergence, term-by-term differentiation); Tome~4, Chapter~10 (values of $\zeta$); Problem~\ref{pb:gaussmeasure}.
\leads Contraction arguments of this kind go back to Kuzmin (1928) and Lévy (1929); the constant $c$ here is far from the
truth, and the notebook estimates the true rate by power iteration. Wirsing (1974) found the exact rate $0.3036630029\ldots$, the second eigenvalue of the transfer
operator on a suitable space. Babenko (1978) described the whole spectrum. Mayer later related this operator to the Selberg zeta
function of the modular group.
\end{solution}

\begin{solution}{pb:levy}
\idea The product $x\,Tx\cdots T^nx$ is both the approximation error of the $n$-th convergent and a product of the values of a
function along the orbit; taking logarithms turns it into a sum, whose terms have, by Gauss--Kuzmin, almost the same law.

\writeup (a) The complete quotient after $a_n$ is $1/T^nx$, so by Tome~1, Proposition~8.1, $x=\frac{p_n/T^nx+p_{n-1}}{q_n/T^nx+q_{n-1}}=\frac{p_n+p_{n-1}T^nx}{q_n+q_{n-1}T^nx}$.
Hence, with $p_nq_{n-1}-p_{n-1}q_n=\pm1$,
\[
\beta_n:=|q_nx-p_n|=\frac{T^nx}{q_n+q_{n-1}T^nx}.
\]
Applied at rank $n-1$, and using $T^{n-1}x=\frac1{a_n+T^nx}$ and $q_n=a_nq_{n-1}+q_{n-2}$, this gives $\beta_{n-1}=\frac1{q_n+q_{n-1}T^nx}$;
so $\beta_n=T^nx\cdot\beta_{n-1}$. With $\beta_0=x$ (as $p_0=0$, $q_0=1$), $\beta_n=x\,Tx\cdots T^nx$, and the formula for $\beta_n$ at rank
$n+1$, $\beta_n=\frac1{q_{n+1}+q_nT^{n+1}x}$, puts it between $\frac1{2q_{n+1}}$ and $\frac1{q_{n+1}}$ (as $q_n\le q_{n+1}$).

(b) $\int_0^1(-\log t)t^m\,dt=\frac1{(m+1)^2}$, and expanding $\frac1{1+t}=\sum_m(-t)^m$ (the partial sums are dominated by $1$),
$\int_0^1\frac{-\log t}{1+t}dt=\sum_{m\ge0}\frac{(-1)^m}{(m+1)^2}=\frac{\pi^2}{12}$.

(c) By (a), $\log q_{n+1}=\sum_{j=0}^n(-\log T^jx)+\theta$ with $-\log2\le\theta\le0$. By Problem~\ref{pb:gausskuzmin}(d), $T^jx$ has a density
$f_j$ with $|f_j(t)-\frac1{(1+t)\log2}|\le c^j$, so
\[
\Bigl|\mathbb E[-\log T^jx]-\frac{\pi^2}{12\log2}\Bigr|\le c^j\int_0^1(-\log t)\,dt=c^j,
\]
using (b). Summing over $j\le n$, $\mathbb E[\log q_{n+1}]=(n+1)\frac{\pi^2}{12\log2}+O\bigl(\frac1{1-c}+\log2\bigr)$.

\tools Tome~1, Chapter~8 (convergents); Tome~4, Chapter~3 (series for $\frac{\pi^2}{12}$); Problem~\ref{pb:gausskuzmin}.
\leads Lévy (1929) proved the almost sure version $\frac1n\log q_n\to\frac{\pi^2}{12\log2}$. For Euclid's algorithm on rationals,
Heilbronn (1969) and Dixon showed that the mean number of division steps for $\frac ab$ with $a<b$ is $\frac{12\log2}{\pi^2}\log b+O(1)$:
the reciprocal of Lévy's constant, as (c) suggests, since each step multiplies the denominator by about $e^{1.1866}$. Porter found
the next term.
\end{solution}

\begin{solution}{pb:benford}
\idea The integers with leading digit $1$ come in blocks $[10^j,2\cdot10^j[$ whose length is a fixed fraction of their position,
so the count never settles. Weighting by $\frac1n$ turns each block into the same contribution $\log2$.

\writeup (a) The number of $n<10^k$ starting with $1$ is $1+10+\dots+10^{k-1}=\frac{10^k-1}9$. At $N=10^k-1$ the proportion is exactly
$\frac19$. At $N=2\cdot10^k-1$ it is $\frac{(10^k-1)/9+10^k}{2\cdot10^k-1}\to\frac{10/9}2=\frac59$. So the proportion has $\liminf\le\frac19$
and $\limsup\ge\frac59$ (in fact these are the extreme values, attained at these $N$ asymptotically), and there is no density.

(b) The block $[d\cdot10^j,(d+1)10^j[$ contributes $\sum\frac1n=\log\frac{d+1}d+O(10^{-j})$ (compare with $\int\frac{dx}x$). For
$10^k\le N<10^{k+1}$ the blocks with $j<k$ contribute $k\log(1+\frac1d)+O(1)$ and the partial block at most $\log10$. Divide by
$\log N=k\log10+O(1)$.

\tools Tome~0, Chapter~5 (comparison of sums and integrals).
\leads Benford's law has been observed in tax records, physical constants and populations, and is used to detect fraud. For sequences
like $2^n$ it holds with natural density, because $n\log_{10}2$ is equidistributed modulo $1$ (Weyl, Chapter~10).
\end{solution}

\begin{solution}{pb:digitsclt}
\idea For $N=2^k$ the bits of a uniform $n$ are independent fair coins. A general $N$ is a union of a few dyadic blocks, each a shifted
copy of the first case, and almost all the mass sits in the big blocks.

\writeup (a) $n\mapsto$ its $k$ bits is a bijection from $\{0,\dots,2^k-1\}$ to $\{0,1\}^k$, so the bits are independent and fair, and
$s(n)$ is their sum.

(b) Write $N=2^{k_1}+\dots+2^{k_r}$ with $k_1>\dots>k_r$, and $M_i=2^{k_1}+\dots+2^{k_{i-1}}$. The interval $[0,N[$ is the disjoint union of
the blocks $B_i=[M_i,M_i+2^{k_i}[$, and for $n\in B_i$, $s(n)=(i-1)+s(n-M_i)$, where $n-M_i$ is uniform on $[0,2^{k_i}[$ when $n$ is uniform
on $B_i$. So the law of $s(n)$ is the mixture, with weights $w_i=2^{k_i}/N$, of the laws of $(i-1)+\mathrm{Bin}(k_i,\frac12)$.

Let $k=k_1$, so $\log_2N\in[k,k+1[$, and fix $h\ge1$. The blocks with $i>h$ have total weight at most $2^{k-h+1}/N\le2^{1-h}$. For
$i\le h$, $k_i\in[k-h,k]$ and the shift is at most $h$; by the de Moivre--Laplace theorem (Tome~3, Chapter~6),
$\bigl((i-1)+\mathrm{Bin}(k_i,\frac12)-\frac12\log_2N\bigr)/\sqrt{\frac14\log_2N}\to N(0,1)$ in law, uniformly in $i\le h$, because the shift
and the change of variance are $O(h)$ and $O(h/k)$ against a standard deviation $\sqrt k/2$. Hence for every $x$, the distribution
function of the normalised $s(n)$ at $x$ is within $2^{1-h}+o(1)$ of $\Phi(x)$. Let $N\to\infty$, then $h\to\infty$.

\tools Tome~3, Chapter~6 (central limit theorem); Tome~3, Chapter~2 (binomial law).
\leads The mean $\frac1N\sum_{n<N}s(n)=\frac12\log_2N+F(\log_2N)$ with $F$ a continuous, nowhere differentiable periodic function
(Trollope 1968, Delange 1975): the blocks make the mean oscillate. Gelfond (1968) proved that $s(n)\bmod2$ is equidistributed in
arithmetic progressions, and Mauduit and Rivat (2010) that it is equidistributed along the primes.
\end{solution}

\begin{solution}{pb:squaresmodpk}
\idea A non-zero square modulo $p^k$ has an even valuation $2j<k$ and a unit part that is a square modulo $p^{k-2j}$. Counting
units that are squares is a matter of group theory: the squaring map on the cyclic unit group has kernel $\pm1$.

\writeup (a) Let $p$ be odd. If $x=y^2\not\equiv0\pmod{p^k}$, then $v_p(x)=2v_p(y)=2j<k$ and $x=p^{2j}u$ with $u$ a unit modulo $p^{k-2j}$ which is a
square modulo $p^{k-2j}$. Conversely every such $p^{2j}u$ is a square. The residue $x$ modulo $p^k$ determines $j$ and $u$ modulo $p^{k-2j}$,
and the units modulo $p^m$ form a cyclic group of even order $\varphi(p^m)$ (Tome~1, Chapter~5), in which the squares are a subgroup of
index $2$. So there are $\frac12\varphi(p^{k-2j})=\frac{p-1}2p^{k-2j-1}$ non-zero squares of valuation $2j$, and adding $0$ gives the formula.
Dividing by $p^k$, $\frac{S(p^k)}{p^k}=p^{-k}+\frac{p-1}{2p}\sum_{0\le j<k/2}p^{-2j}\to\frac{p-1}{2p}\cdot\frac1{1-p^{-2}}=\frac p{2(p+1)}$.

(b) For $p=2$ the units modulo $2^m$ that are squares are those $\equiv1\pmod8$ when $m\ge3$ (and only $1$ for $m\le2$): a proportion
$\frac14$ of the $2^{m-1}$ units, that is $2^{m-3}$. So $\frac{S(2^k)}{2^k}=\sum_{2j\le k-3}2^{-2j}\cdot\frac18+O(2^{-k/2})\to\frac18\cdot\frac43=\frac16$.

(c) By the Chinese remainder theorem, $x$ is a square modulo $mn$ (coprime) if and only if it is a square modulo $m$ and modulo $n$,
so $S(mn)=S(m)S(n)$. For $m=\prod_{p\le y}p$, the proportion of squares is $\prod_{p\le y}\frac{p+1}{2p}\le\prod_{3\le p\le y}\frac23\to0$. An
integer $a$ that is a square modulo every $m$ is a square: otherwise write $a=b^2c$ with $c\ne1$ squarefree (or $a<0$), and choose,
by Chinese remainders and quadratic reciprocity, a prime $q\nmid2a$ with $\bigl(\frac cq\bigr)=-1$; then $a$ is not a square modulo $q$.

\tools Tome~1, Chapters~5, 6 and~12 (cyclic unit groups, reciprocity, $p$-adic squares); Tome~2, Chapter~2.
\leads In $p$-adic language, the squares of $\Z_p$ have Haar measure $\frac p{2(p+1)}$ ($\frac16$ for $p=2$). The last statement of (c)
is the simplest case of a \emph{local--global principle}: a property that holds modulo every $m$ is forced
to hold over $\Z$. Such principles are a recurring theme of Part~III.
\end{solution}

\begin{solution}{pb:hashing}
\idea Two distinct points determine a line, so the values of a random line $ax+b$ at two distinct points are an arbitrary pair.
Reducing modulo $m$ loses a little uniformity, and that is all.

\writeup (a) For $x\ne y$ in $\Z/p\Z$ the linear map $(a,b)\mapsto(ax+b,ay+b)$ has determinant $x-y\ne0$, so it is a bijection of
$(\Z/p\Z)^2$, and maps the uniform law to the uniform law.

(b) By (a), $P(h(x)=h(y))$ is the proportion of pairs $(r,s)\in\{0,\dots,p-1\}^2$ with $r\equiv s\pmod m$. For each $r$ there are at most
$\lceil p/m\rceil\le\frac{p-1}m+1$ values of $s$ in its class, so the probability is at most $\frac1p\bigl(\frac{p-1}m+1\bigr)\le\frac1m+\frac1p$.

(c) The expected number of colliding pairs among $n$ keys is at most $\binom n2(\frac1{n^2}+\frac1p)\le\frac12+\frac{n^2}{2p}$. By Markov's
inequality the probability of at least one collision is at most that, so no collision happens with probability at least
$\frac12-\frac{n^2}{2p}$.

\tools Tome~0, Chapter~7 (linear maps); Tome~3, Chapters~2 and~4 (linearity of expectation, Markov's inequality).
\leads (c) is the first level of the Fredman--Komlós--Szemerédi perfect hashing scheme (1984), which stores $n$ keys in $O(n)$ space
with worst-case constant lookup time. Pairwise independence is enough here, which is why two random numbers $a,b$ suffice instead
of a fully random function.
\end{solution}

\begin{solution}{pb:cdg}
\idea The walk is a random binary-like expansion $X_n=\sum2^{n-1-j}\varepsilon_j$, so its Fourier transform is a product. A factor is
small as soon as the doubling orbit of $t/p$ visits the far half of the circle, and an orbit cannot stay near $0$ for more than
$\log_2p$ steps, because doubling pushes it away.

\writeup (a) $X_n$ takes at most $3^n$ values, so the law of $X_n$ puts mass $1$ on a set $A$ with $|A|\le3^n$, which has uniform measure
$|A|/p$. The total variation distance is at least $1-\frac{3^n}p$.

(b) By induction, $X_n\equiv\sum_{j<n}2^{n-1-j}\varepsilon_j$. The $\varepsilon_j$ are independent and
$\mathbb E[e^{2i\pi s\varepsilon/p}]=\frac{1+2\cos(2\pi s/p)}3$, so the transform is $\prod_{j<n}\frac{1+2\cos(2\pi2^{n-1-j}t/p)}3$, which is the
stated product after $j\mapsto n-1-j$.

(c) Let $\theta_j\in\,]-\frac12,\frac12]$ be the representative of $2^jt/p$ modulo $1$; $\theta_j\ne0$ is a multiple of $\frac1p$, so
$|\theta_j|\ge\frac1p$. If $|\theta_j|<\frac14$, then $2\theta_j\in\,]-\frac12,\frac12[$ and $\theta_{j+1}=2\theta_j$. If
$|\theta_j|<\frac14$ for all $j_0\le j<j_0+L$, then $|\theta_{j_0+L-1}|=2^{L-1}|\theta_{j_0}|\ge\frac{2^{L-1}}p\ge1$, a contradiction. So some
$j$ in the window has $|\theta_j|\ge\frac14$, that is $\cos(2\pi\theta_j)\le0$.

(d) When $\cos\le0$, $\bigl|\frac{1+2\cos}3\bigr|\le\frac13$; every factor has modulus at most $1$. Among $j<n$ there are
$\lfloor n/L\rfloor$ disjoint windows of length $L$, so $|\hat P(t)|\le3^{-\lfloor n/L\rfloor}$ for $t\ne0$, where $\hat P$ is the transform of
(b). By Cauchy--Schwarz and Plancherel on $\Z/p\Z$ (Tome~4, Chapter~7),
\[
d_n=\frac12\sum_x\Bigl|P(X_n=x)-\frac1p\Bigr|\le\frac12\sqrt p\Bigl(\sum_x\Bigl|P(X_n=x)-\frac1p\Bigr|^2\Bigr)^{1/2}=\frac12\Bigl(\sum_{t\ne0}|\hat P(t)|^2\Bigr)^{1/2}\le\frac12\sqrt p\,3^{-\lfloor n/L\rfloor}.
\]
So $d_n\le\varepsilon$ once $\lfloor n/L\rfloor\ge\log_3\frac{\sqrt p}{2\varepsilon}$, which is $n=O(\log^2p)$ for fixed $\varepsilon$.

\tools Tome~3, Chapter~6 (characteristic functions); Tome~4, Chapter~7 (Fourier analysis on $\Z/p\Z$); Tome~3, Chapter~8 (total
variation).
\leads Chung, Diaconis and Graham proved that $O(\log p\log\log p)$ steps always suffice and that about $1.02\log_2p$ steps suffice for almost all
$p$, and that for some $p$ (of the form $2^k-1$) the order $\log p\log\log p$ is needed. Compare with the walk $X_{n+1}=X_n+\varepsilon_n$,
which needs about $p^2$ steps (Chapter~11): multiplying by $2$ spreads the walk out exponentially fast. Hildebrand and Eberhard
and Varjú have since sharpened the picture.
\end{solution}

\section*{Map of the chapter}

\begin{center}
\small
\adjustbox{max width=\linewidth}{\begin{tabular}{lll}
Problem & Topic & Tomes used\\\hline
\ref{pb:nouniform} & no uniform integer, densities & 3\\
\ref{pb:valuation} & valuations of a random integer & 1, 2\\
\ref{pb:consecsqfree} & consecutive squarefree numbers & 2, 4\\
\ref{pb:pairwise} & pairwise coprime tuples & 2, 3, 5\\
\ref{pb:parity} & a random fraction in lowest terms & 1\\
\ref{pb:nsq1sqfree} & squarefree values of $n^2+1$ & 1, 5\\
\ref{pb:twosqfree} & sums of two squarefree numbers & 1, 2\\
\ref{pb:phiovern} & the law of $\varphi(n)/n$ & 1, 3, 5\\
\ref{pb:gcdmean} & law, mean and log-mean of the gcd & 1, 3, 5\\
\ref{pb:fibcoprime} & coprime Fibonacci numbers & 0\\
\ref{pb:visible} & visible lattice points & 1, 5\\
\ref{pb:orchard} & Pólya's orchard, Minkowski & 0, 1\\
\ref{pb:farey} & Farey fractions & 1, 5\\
\ref{pb:euclidgame} & Euclid's game & 0, 1\\
\ref{pb:gaussmeasure} & the Gauss measure & 1, 3\\
\ref{pb:gausskuzmin} & the Gauss--Kuzmin theorem & 0, 4\\
\ref{pb:levy} & Lévy's constant in mean & 1, 4\\
\ref{pb:benford} & leading digits, Benford & 0\\
\ref{pb:digitsclt} & binary digits, CLT & 3\\
\ref{pb:squaresmodpk} & squares modulo $p^k$ & 1, 2\\
\ref{pb:hashing} & universal hashing & 0, 3\\
\ref{pb:cdg} & the walk $X\mapsto2X\pm1$ modulo $p$ & 3, 4\\
\end{tabular}}
\end{center}

\section*{Sources and further reading}

\begin{itemize}
\item P.-S.~Laplace, \emph{Essai philosophique sur les probabilités}, 1814; in \emph{Œuvres complètes}, vol.~VII, Gauthier-Villars,
p.~cliii (the epigraph: ``On voit, par cet Essai, que la théorie des probabilités n'est, au fond, que le bon sens réduit au
calcul'', our translation).
\item M.~Iosifescu and C.~Kraaikamp, \emph{Metrical Theory of Continued Fractions}, Kluwer, 2002, Chapter~2 (Gauss's diary entry of
25 October 1800, his letter to Laplace of 30 January 1812, and the work of Kuzmin, Lévy, Szüsz and Wirsing); D.~Lascu, On the Szüsz's
solution to Gauss' problem, arXiv:1010.4432.
\item R.~O.~Kuzmin, On a problem of Gauss, \emph{Doklady Akademii Nauk SSSR A} (1928), 375--380; P.~Lévy, Sur les lois de probabilité
dont dépendent les quotients complets et incomplets d'une fraction continue, \emph{Bulletin de la Société Mathématique de France} 57
(1929), 178--194; E.~Wirsing, On the theorem of Gauss--Kusmin--Lévy and a Frobenius-type theorem for function spaces, \emph{Acta
Arithmetica} 24 (1974), 507--528.
\item A.~Ya.~Khinchin, \emph{Continued Fractions}, University of Chicago Press, 1964 (Dover reprint 1997).
\item M.~Kac, \emph{Statistical Independence in Probability, Analysis and Number Theory}, Carus Mathematical Monographs 12, 1959.
\item G.~Tenenbaum, \emph{Introduction to Analytic and Probabilistic Number Theory}, 3rd ed., AMS, 2015 (densities, $\varphi(n)/n$,
squarefree values).
\item T.~Estermann, Einige Sätze über quadratfreie Zahlen, \emph{Mathematische Annalen} 105 (1931), 653--662.
\item G.~Pólya, Zahlentheoretisches und Wahrscheinlichkeitstheoretisches über die Sichtweite im Walde, \emph{Archiv der Mathematik und
Physik} 27 (1918), 135--142; G.~Pólya and G.~Szegő, \emph{Problems and Theorems in Analysis II}, Springer, 1976.
\item A.~J.~Cole and A.~J.~T.~Davie, A game based on the Euclidean algorithm and a winning strategy for it, \emph{Mathematical
Gazette} 53 (1969), 354--357.
\item G.~H.~Hardy and E.~M.~Wright, \emph{An Introduction to the Theory of Numbers}, 6th ed., Oxford, 2008, Chapter~3 (Farey
fractions).
\item F.~P.~Boca, C.~Cobeli and A.~Zaharescu, Distribution of lattice points visible from the origin, \emph{Communications in
Mathematical Physics} 213 (2000), 433--470.
\item J.~L.~Carter and M.~N.~Wegman, Universal classes of hash functions, \emph{Journal of Computer and System Sciences} 18 (1979),
143--154; M.~L.~Fredman, J.~Komlós and E.~Szemerédi, Storing a sparse table with $O(1)$ worst case access time, \emph{Journal of the
ACM} 31 (1984), 538--544.
\item F.~R.~K.~Chung, P.~Diaconis and R.~L.~Graham, Random walks arising in random number generation, \emph{Annals of Probability} 15
(1987), 1148--1165.
\item C.~Mauduit and J.~Rivat, Sur un problème de Gelfond: la somme des chiffres des nombres premiers, \emph{Annals of Mathematics}
171 (2010), 1591--1646.
\end{itemize}
